% =============================================================================
% БИЛЕТ 10
% =============================================================================

\ticket{10}{}
\question{Теорема о диффеоморфизме. Теорема о локальном диффеоморфизме.
Продолжение до диффеоморфизма}

% ═══════════════════════════════════════════════════════════════════════════
%  ОПРЕДЕЛЕНИЕ ДИФФЕОМОРФИЗМА
% ════════════════════════════════════════════════════════════════════════════════════

\begin{defn}[Диффеоморфизм]
  Пусть $U, V \subset \mathbb{R}^m$ --- непустые открытые множества.
  Отображение $F : U \to V$ называется \emph{$C^1$-диффеоморфизмом}, если:
  \begin{enumerate}
    \item $F$ --- биекция $U$ на $V$;
    \item $F \in C^1(U, V)$;
    \item $F^{-1} \in C^1(V, U)$.
  \end{enumerate}
\end{defn}

\begin{note}
  Пункт~3 необходимо требовать отдельно: из гладкости $F$ гладкость $F^{-1}$
  не следует автоматически. Контрпример: $f(x) = x^3$ --- биекция,
  $f \in C^1(\mathbb{R})$, но $(f^{-1})'(0) = +\infty$, так как $f'(0) = 0$.
  Именно поэтому в теоремах ниже требуется $\mathcal{J}F \neq 0$: нулевой
  якобиан означает «схлопывание» отображения, что ломает гладкость обратного.
\end{note}

% ═══════════════════════════════════════════════════════════════════════════
%  ВСПОМОГАТЕЛЬНЫЕ ЛЕММЫ
% ═══════════════════════════════════════════════════════════════════════════

\begin{lemma}[Неравенство Лагранжа 2.0]
  \label{lemma:lagrange2}
  Пусть $B_\delta(x^0) \subset \mathbb{R}^n$,
  $F \in \mathrm{DIF}(B_\delta(x^0), \mathbb{R}^n)$.
  Тогда для любых $x, y \in B_\delta(x^0)$ существует $\theta \in (0,1)$ такое, что:
  \[
    \|F(x) - F(y)\|
    \;\leq\;
    \|\mathcal{D}F\bigl(x + \theta(y-x)\bigr)\| \cdot \|y - x\|.
  \]
\end{lemma}

\begin{proof}
  Так как шар $B_\delta(x^0)$ выпуклый, отрезок
  $\{x + t(y-x),\; t \in [0,1]\}$ целиком лежит в $B_\delta(x^0)$.

  Рассмотрим вектор-функцию одного переменного:
  \[
    H(t) \;:=\; F\bigl(x + t(y-x)\bigr), \quad t \in [0,1].
  \]
  Она дифференцируема как композиция дифференцируемых отображений.
  По теореме о дифференцировании композиции:
  \[
    H'(t) = \mathcal{D}F\bigl(x + t(y-x)\bigr) \cdot (y - x).
  \]
  По теореме Лагранжа о среднем для вектор-функций:
  \[
    \|H(1) - H(0)\| \;\leq\; \|H'(\theta)\| \cdot 1
    \;=\; \|\mathcal{D}F\bigl(x + \theta(y-x)\bigr) \cdot (y-x)\|
    \;\leq\; \|\mathcal{D}F\bigl(x + \theta(y-x)\bigr)\| \cdot \|y-x\|.
  \]
  Остаётся заметить, что $H(1) = F(y)$ и $H(0) = F(x)$.
\end{proof}

\begin{lemma}[Неравенство Лагранжа 3.0]
  \label{lemma:lagrange3}
  Пусть $B_\delta(x^0) \subset \mathbb{R}^n$,
  $F \in C^1(B_\delta(x^0), \mathbb{R}^n)$.
  Тогда для любых $x, y \in B_\delta(x^0)$:
  \[
    \bigl\|F(y) - F(x) - \mathcal{D}F(x^0)(y-x)\bigr\|
    \;\leq\;
    \sup_{z \in [x,y]}\bigl\|\mathcal{D}F(z) - \mathcal{D}F(x^0)\bigr\|
    \cdot \|y - x\|.
  \]
  \textbf{Смысл:} правая часть --- это ошибка линейного приближения $F$
  матрицей $\mathcal{D}F(x^0)$. Если мы в маленьком шаре (где $\mathcal{D}F$
  почти не меняется), то эта ошибка мала.
\end{lemma}

\begin{proof}
  Введём вспомогательное отображение:
  \[
    L(z) \;:=\; F(z) - \mathcal{D}F(x^0) \cdot z.
  \]
  Оно $C^1$-гладкое как разность $C^1$-отображения и линейного.
  Его матрица Якоби:
  \[
    \mathcal{D}L(z) = \mathcal{D}F(z) - \mathcal{D}F(x^0).
  \]
  Применяем лемму~\ref{lemma:lagrange2} к $L$:
  \[
    \|L(y) - L(x)\|
    \;\leq\;
    \sup_{z \in [x,y]} \|\mathcal{D}L(z)\| \cdot \|y - x\|
    \;=\;
    \sup_{z \in [x,y]} \|\mathcal{D}F(z) - \mathcal{D}F(x^0)\| \cdot \|y - x\|.
  \]
  Осталось раскрыть $L(y) - L(x)$:
  \[
    L(y) - L(x)
    = \bigl(F(y) - \mathcal{D}F(x^0) \cdot y\bigr)
      - \bigl(F(x) - \mathcal{D}F(x^0) \cdot x\bigr)
    = F(y) - F(x) - \mathcal{D}F(x^0)(y-x). \qed
  \]
\end{proof}

% ═══════════════════════════════════════════════════════════════════════════
%  ТЕОРЕМА О ДИФФЕОМОРФИЗМЕ (ГЛОБАЛЬНАЯ)
% ═══════════════════════════════════════════════════════════════════════════

\begin{theorem}[О диффеоморфизме]
  \label{theorem:global}
  Пусть $\Omega \subset \mathbb{R}^n$ открыто,
  $F \in C^1(\Omega, \mathbb{R}^n)$,
  $F$ инъективна на $\Omega$,
  $\mathcal{J}F(x) \neq 0$ для всех $x \in \Omega$.
  Тогда $F$ является $C^1$-диффеоморфизмом $\Omega$ на $F(\Omega)$.
\end{theorem}

\begin{note}
  Пункты 1--2 определения диффеоморфизма даны прямо по условию:
  биективность на $F(\Omega)$ --- по условию и по определению образа,
  $F \in C^1$ --- по условию.
  Нетривиален только пункт~3: нужно доказать, что $F^{-1} \in C^1$.
\end{note}

\begin{proof}
  \textbf{Шаг 1. Непрерывность $F^{-1}$.}

  По теореме об открытом отображении (применима, так как $F \in C^1$
  и $\mathcal{J}F \neq 0$): для любого открытого $U \subset \Omega$
  множество $F(U)$ открыто.

  Используем критерий непрерывности: отображение непрерывно тогда и
  только тогда, когда прообраз любого открытого открыт.
  Для $F^{-1}$ и произвольного открытого $U \subset \Omega$:
  \[
    (F^{-1})^{-1}(U) \;=\; F(U) \;\text{--- открыто.}
  \]
  Следовательно, $F^{-1}$ \textbf{непрерывна}.

  \medskip
  \textbf{Шаг 2. Формула для производной $F^{-1}$.}

  Так как $\mathcal{J}F(x) \neq 0$, матрица $\mathcal{D}F(x)$ обратима.
  По теореме о производной обратного отображения:
  \[
    \mathcal{D}F^{-1}(y)
    \;=\; \bigl(\mathcal{D}F\bigl(F^{-1}(y)\bigr)\bigr)^{-1}.
  \]

  \medskip
  \textbf{Шаг 3. Непрерывность $\mathcal{D}F^{-1}$.}

  Запишем $y \mapsto \mathcal{D}F^{-1}(y)$ как цепочку трёх отображений:
  \[
    y
    \;\xrightarrow{A}\; F^{-1}(y) = x
    \;\xrightarrow{B}\; \mathcal{D}F(x)
    \;\xrightarrow{C}\; \bigl(\mathcal{D}F(x)\bigr)^{-1}.
  \]
  \begin{itemize}
    \item $A$ непрерывно: доказано в шаге~1.
    \item $B$ непрерывно: $F \in C^1$, значит $\mathcal{D}F$ непрерывна.
    \item $C$ непрерывно: элементы $M^{-1}$ --- рациональные функции
          элементов $M$ (знаменатель $\det M \neq 0$ на $GL(n)$),
          значит взятие обратной матрицы непрерывно.
  \end{itemize}
  Композиция непрерывных непрерывна, значит $F^{-1} \in C^1$. \qed
\end{proof}

% ═══════════════════════════════════════════════════════════════════════════
%  ТЕОРЕМА О ЛОКАЛЬНОМ ДИФФЕОМОРФИЗМЕ
% ═══════════════════════════════════════════════════════════════════════════

\begin{theorem}[О локальном диффеоморфизме]
  \label{theorem:local}
  Пусть $F \in C^1(\Omega, \mathbb{R}^n)$ и
  $\mathcal{J}F(x) \neq 0$ для всех $x \in \Omega$.
  Тогда для каждой точки $x^0 \in \Omega$ существует шар
  $B_\delta(x^0) \subset \Omega$ такой, что
  $F\big|_{B_\delta(x^0)}$ --- $C^1$-диффеоморфизм.
\end{theorem}

\begin{note}
  Почему только локально? Якобиан может быть ненулевым всюду, но глобальной
  инъективности при этом нет. Пример: полярные координаты
  $F(r,\varphi) = (r\cos\varphi,\; r\sin\varphi)$
  имеют $\mathcal{J}F = r > 0$ при $r > 0$, но
  $(r, \varphi)$ и $(r, \varphi + 2\pi)$ дают одну точку.
\end{note}

\begin{proof}
  По теореме~\ref{thm:global} достаточно найти шар $B_\delta(x^0)$,
  на котором $F$ \textbf{инъективна}: условия $F \in C^1$ и
  $\mathcal{J}F \neq 0$ уже даны по условию теоремы.

  \medskip
  \textbf{Шаг 1. Оценка снизу для обратимого оператора.}

  Обозначим $A := \mathcal{D}F(x^0)$. Так как $\mathcal{J}F(x^0) \neq 0$,
  оператор $A$ обратим. По основной геометрической теореме:
  \[
    \exists\, c > 0 \;:\quad \|Ah\| \;\geq\; c\,\|h\|
    \quad \forall h \in \mathbb{R}^n. \tag{$*$}
  \]
  Смысл: обратимый оператор не может «схлопнуть» ненулевой вектор ---
  он всегда растягивает хотя бы на константу $c$.

  \medskip
  \textbf{Шаг 2. Выбор шара, где $\mathcal{D}F$ почти постоянна.}

  Так как $F \in C^1$, матрица $\mathcal{D}F$ непрерывна.
  Значит в точке $x^0$ существует $\delta > 0$ такое, что:
  \[
    \bigl\|\mathcal{D}F(z) - \mathcal{D}F(x^0)\bigr\| < \frac{c}{2}
    \quad \forall z \in B_\delta(x^0). \tag{$\vee$}
  \]
  Смысл: в достаточно маленьком шаре $\mathcal{D}F$ почти не отличается
  от $A = \mathcal{D}F(x^0)$, то есть ошибка линейного приближения мала.

  \medskip
  \textbf{Шаг 3. Инъективность на $B_\delta(x^0)$.}

  Для произвольных $x, y \in B_\delta(x^0)$ добавим и вычтем $A(y-x)$:
  \begin{align*}
    \|F(y) - F(x)\|
    &= \|F(y) - F(x) - A(y-x) + A(y-x)\| \\[4pt]
    &\geq \|A(y-x)\| - \|F(y) - F(x) - A(y-x)\|
      \tag{нер-во треугольника} \\[4pt]
    &\geq c\,\|y-x\|
      \;-\; \sup_{z\in[x,y]}\|\mathcal{D}F(z)-\mathcal{D}F(x^0)\|\cdot\|y-x\|
      \tag{$(*)$ и лемма~\ref{lemma:lagrange3}} \\[4pt]
    &\geq c\,\|y-x\| - \frac{c}{2}\,\|y-x\|
      \tag{из $(\vee)$} \\[4pt]
    &= \frac{c}{2}\,\|y-x\|.
  \end{align*}
  Если $x \neq y$, то $\|F(y) - F(x)\| \geq \tfrac{c}{2}\|y-x\| > 0$,
  значит $F(y) \neq F(x)$ --- $F$ \textbf{инъективна} на $B_\delta(x^0)$.

  По теореме~\ref{theorem:global} заключаем, что $F$ ---
  $C^1$-диффеоморфизм $B_\delta(x^0)$ на $F(B_\delta(x^0))$. \qed
\end{proof}

% ═══════════════════════════════════════════════════════════════════════════
%  ПРОДОЛЖЕНИЕ ДО ДИФФЕОМОРФИЗМА
% ═══════════════════════════════════════════════════════════════════════════

\begin{theorem}[Продолжение до диффеоморфизма]
  \label{theorem:extension}
  Пусть $1 \leq k < m$, $\Omega \subset \mathbb{R}^k$ открыто,
  $\Phi \in C^1(\Omega, \mathbb{R}^m)$ и
  $\mathrm{rg}\,\mathcal{D}\Phi(x) = k$ для всех $x \in \Omega$.

  Отождествим $\mathbb{R}^k \equiv \mathbb{R}^k \times \{0\} \subset \mathbb{R}^m$.
  Тогда для любой точки $x^0 \in \Omega$ существуют
  $m$-мерная окрестность $U(x^0, 0) \subset \mathbb{R}^m$ и
  отображение $\overline{\Phi} \in C^1(U(x^0,0),\, \mathbb{R}^m)$ такие, что:
  \begin{enumerate}
    \item $\overline{\Phi}\big|_{U(x^0,0)\,\cap\,\mathbb{R}^k} = \Phi$;
    \item $\overline{\Phi}$ является $C^1$-диффеоморфизмом $U(x^0,0)$
          на открытое множество $V \subset \mathbb{R}^m$.
  \end{enumerate}
\end{theorem}

\begin{note}
  Смысл: $\Phi$ отображает $k$-мерную область в $m$-мерное пространство
  ($k < m$). Само по себе оно не может быть диффеоморфизмом
  (размерности не совпадают). Но мы можем \emph{доопределить}
  $\Phi$ до квадратного отображения $\overline{\Phi}$ на $m$-мерной
  окрестности, и тогда $\overline{\Phi}$ уже диффеоморфизм.
\end{note}

\begin{proof}
  \textbf{Шаг 1. Конструкция расширенного отображения.}

  Определим $\overline{\Phi} : \Omega \times \mathbb{R}^{m-k} \to \mathbb{R}^m$:
  \[
    \overline{\Phi}(x,\, y) \;:=\; \Phi(x) + (0_k,\, y),
    \quad x \in \Omega,\; y \in \mathbb{R}^{m-k}.
  \]
  При $y = 0$ получаем $\overline{\Phi}(x, 0) = \Phi(x)$, то есть
  $\overline{\Phi}$ совпадает с $\Phi$ на $\mathbb{R}^k$ --- пункт~1 выполнен.

  \medskip
  \textbf{Шаг 2. Якобиан $\overline{\Phi}$ ненулевой.}

  Матрица Якоби $\overline{\Phi}$ в точке $(x^0, 0)$ блочно-треугольная:
  \[
    \mathcal{D}\overline{\Phi}(x^0, 0)
    \;=\;
    \begin{pmatrix}
      \mathcal{D}\Phi(x^0) & 0 \\[4pt]
      0 & E_{(m-k)\times(m-k)}
    \end{pmatrix}.
  \]
  Её определитель:
  \[
    \mathcal{J}\overline{\Phi}(x^0, 0)
    \;=\; \det\mathcal{D}\Phi(x^0) \cdot \underbrace{\det E_{m-k}}_{=\,1}
    \;\neq\; 0,
  \]
  так как $\mathrm{rg}\,\mathcal{D}\Phi(x^0) = k$ означает
  $\det\mathcal{D}\Phi(x^0) \neq 0$.

  Из непрерывности якобиана по принципу локализации:
  $\mathcal{J}\overline{\Phi} \neq 0$ в некотором шаре вокруг $(x^0, 0)$.

  \medskip
  \textbf{Шаг 3. Применение теоремы о локальном диффеоморфизме.}

  Отображение $\overline{\Phi} \in C^1$ и $\mathcal{J}\overline{\Phi} \neq 0$
  в шаре вокруг $(x^0, 0)$. По теореме~\ref{thm:local} существует
  окрестность $U(x^0, 0)$ такая, что $\overline{\Phi}\big|_{U(x^0,0)}$
  --- $C^1$-диффеоморфизм на $V := \overline{\Phi}(U(x^0,0))$ --- пункт~2
  выполнен. \qed
\end{proof}

\question{Замена переменной и интегрирование по частям в определённом интеграле}

% Ответ...

